\documentclass[11pt]{article}
\usepackage[T1]{fontenc}
\usepackage{lmodern,amsmath,amssymb,amsthm}
\usepackage[margin=1in]{geometry}
\usepackage[hidelinks]{hyperref}
\newtheorem{theorem}{Theorem}
\newcommand{\Q}{\mathbb Q}
\newcommand{\R}{\mathbb R}
\newcommand{\Z}{\mathbb Z}
\title{A Certified Four-Zero Disk for the Elliptic Curve 234446.a1}
\author{\small\href{mailto:admin@cosmosgeek.org}{admin@cosmosgeek.org}}
\date{Computational working paper}
\begin{document}
\maketitle
\begin{abstract}
We certify that the Hasse--Weil \(L\)-function of the elliptic curve
\(y^2+xy=x^3-x^2-79x+289\) has exactly four zeros, counted with
multiplicity, within distance \(10^{-57}\) of \(1\).
An interval computation encloses its second and fourth completed
central derivatives. A strip bound for trapezoidal quadrature and a
Taylor remainder bound turn these enclosures into a zero count.
The count leaves open whether the central point itself has order two
or four.
\end{abstract}

\section{Statement and normalization}
The curve \(E\) in the abstract has conductor \(N=234446=2\cdot117223\)
and Mordell--Weil rank four \cite{curve}. Its two finite local root
numbers are \(+1,-1\) at \(2,117223\), respectively. Including the real
place gives sign \(+1\). Write
\[
 c=\frac{2\pi}{\sqrt N},\qquad
 \Lambda(s)=c^{-s}\Gamma(s)L(E,s),\qquad
 f(z)=\sum_{n\ge1}a_n e^{2\pi i n z}.
\]
Here \(f\) is the normalized modular newform of \(E\). The theorem
that analytic rank zero implies Mordell--Weil rank zero gives
\(L(E,1)=0\) \cite{Wiles}. The even functional equation then gives
\(\Lambda'(1)=L'(E,1)=0\).

\begin{theorem}
There are exactly four zeros of \(L(E,s)\), counting multiplicity,
in \(|s-1|<10^{-57}\). If the central zero has order two, the other
two zeros in that disk are distinct and are either real or on the
line \(\Re s=1\).
\end{theorem}
The LMFDB reliability statement discusses why numerical proximity to
zero does not certify analytic rank four for this curve \cite{reliability}.
The specific contribution here is a reproducible disk count with
explicit truncation and discretization errors.

\section{Enclosing the two moments}
Let \(g(t)=e^t f(ie^t/\sqrt N)\). The Mellin integral and Fricke
symmetry give
\[
 \Lambda(1+u)=\int_{\R}e^{ut}g(t)\,dt,\qquad g(-t)=g(t).
\]
Both identities extend by analyticity on their respective domains.
For \(j=2,4\) let \(I_j=\int_\R t^jg(t)\,dt=\Lambda^{(j)}(1)\).
Set \(h=1/32,\ d=7/5\). Approximate \(I_j\) by the finite sum
\[
 Q_j=2h\sum_{i=1}^{320}(ih)^j e^{ih}
       \sum_{n=1}^{M_i}a_n \exp(-nc e^{ih}).
\]
Here \(ih=i/32\) is real, and \(i\) is an integer summation index.
Choose \(M_i\) by rounding \(300/(ce^{ih})\) up, using the lower
endpoint of an enclosure of \(ce^{ih}\) in the denominator. With the
enclosures below, \(M_i\le25001\).

\paragraph{Infinite mesh.}
The Fourier series is analytic for \(|\Im t|<\pi/2\). Put
\(\alpha=c\cos d>0.002\). Using \(|a_n|\le d(n)\sqrt n\le2n\)
and \(\sum_{n\ge1}n e^{-nv}=(2\sinh(v/2))^{-2}\le v^{-2}\)
gives
\[
 |g(x\pm id)|\le 2\alpha^{-2} e^{-x}\qquad(x\ge0).
\]
The other half-lines satisfy the reflected bounds. The two boundary
integrals of \(|t^j g(t)|\) therefore sum to at most
\[
 A_2=8(2+d^2)/\alpha^2,\qquad A_4=8(24+4d^2+d^4)/\alpha^2 .
\]
The infinite trapezoidal error is bounded by
\(A_j/(\exp(2\pi d/h)-1)\). This follows by contour shifting the
Fourier transform at the nonzero frequencies \(2\pi k/h\) to the
two boundary lines and summing the geometric series \cite{TW}.
The bounds are \(4.5\cdot10^{-116}\) for \(j=2\) and
\(4.1\cdot10^{-115}\) for \(j=4\).

\paragraph{Two tails.}
Beyond \(x=10\) use
\(|g(x)|\le8e^{x-ce^x}\). The logarithmic derivative of
\(x^je^{x-ce^x}\) is \(j/x+1-ce^x<-280\) there. Thus the omitted
mesh contributes at most
\[
 \frac{16h\,10^j e^{10}\exp(-ce^{10})}{\exp(280h)-1}
 <10^{-118}\qquad(j=2,4).
\]
At mesh point \(i\) put \(q_i=\exp(-ce^{ih})\).
The discarded Fourier coefficients have total absolute contribution
to the inner sum at most
\[
 2q_i^{M_i+1}
 \left(\frac{M_i+1}{1-q_i}+
       \frac{q_i}{(1-q_i)^2}\right).
\]
Using \(q_i^{M_i+1}\le e^{-300}\),
\(q_i\le e^{-0.012}\), and summing the outer weights by
\(20\cdot 10^j e^{10}\) gives bounds \(10^{-116}\) and \(10^{-114}\),
respectively. In particular
\[
 |I_2-Q_2|<10^{-113},\qquad |I_4-Q_4|<10^{-112}.
\]

\paragraph{Finite arithmetic.}
Every needed prime other than 2 is good. The program below uses the
short model \(Y^2=X^3-1267X+17230\), obtained from
\(X=4x-1,\ Y=8y+4x\), for the odd-prime point counts. The bad
coefficient at 2 is \(-1\). The Euler factors supply every
\(a_n\) through Hecke recursion. The Machin identity
\(\pi=16\arctan(1/5)-4\arctan(1/239)\) with its alternating series
and an integer square root give interval endpoints for \(c\).
The interval endpoints in the Horner sums are fixed-point integers
with scale \(10^{150}\). Python's Decimal exponential is
correctly rounded \cite{decimal}. The neighboring decimal numbers
and outward integer division bound every exponential.
Executing the program gives the deliberately widened enclosures
\[
 -1.207\cdot10^{-123}<Q_2<-1.206\cdot10^{-123},
 \qquad 16541.5867<Q_4<16541.5869.
\]
Hence \(|\Lambda''(1)|<2\cdot10^{-113}\) and \(\Lambda^{(4)}(1)>16541\).
The small negative interval above describes the finite sum only.

\section{The disk count}
Let \(A(u)=\Lambda(1+u)\). For \(|u|\le1\) the Mellin
integral, reflection, and the coefficient bound give
\[
 |A(u)|\le4\sum_{n\ge1}n\int_1^\infty t e^{-cnt}\,dt
 =4\left(\frac1{c(e^c-1)}
  -\frac{\log(1-e^{-c})}{c^2}\right)<2\cdot10^5 .
\]
The last inequality already follows from \(c>0.0129\).
Cauchy's bound at radius one controls the Taylor remainder from
degree six onward by \(2\cdot10^5 R^6/(1-R)\) on \(|u|=R\).
The constant term is zero, and the odd coefficients vanish.
For \(R=10^{-57}\),
\[
 \tfrac12 |I_2|R^2+\frac{2\cdot10^5 R^6}{1-R}
 <\frac{16541}{24}R^4.
\]
Rouch\'e's theorem, applied to \(I_4u^4/24\), proves the count.
Completion does not alter the zeros in this disk. If the central
order is two, evenness and conjugation symmetry force the two
remaining zeros into one symmetric real or one symmetric
imaginary pair in the \(u\)-coordinate. Their coincidence would
place them at the origin.

\section*{Reproducing the finite enclosures}
Run the following code with Python 3.8 or later. Values printed are
endpoints before adding the three analytic errors above.
\begingroup\fontsize{8}{9.6}\selectfont
\font\codeface=ec-lmtt10 at 8pt\let\ttfamily\codeface
\begin{verbatim}
"""Certified central moments for Cremona 234446a1. Python 3, standard library."""
from decimal import Decimal, getcontext
from math import isqrt

D = 150
B = 10**D
getcontext().prec = D + 22

def ceildiv(a, b):
    return -((-a)//b)

def mul(x, y):
    products = [u*v for u in x for v in y]
    return min(products)//B, ceildiv(max(products), B)

def decint(x, upper):
    n, d = x.as_integer_ratio()
    return ceildiv(n*B, d) if upper else n*B//d

def exponential(x):
    lo = Decimal(x[0]).scaleb(-D).exp().next_minus()
    hi = Decimal(x[1]).scaleb(-D).exp().next_plus()
    return decint(lo, False), decint(hi, True)

def atan_inv(q):
    lo = hi = j = 0
    while True:
        d = (2*j+1)*q**(2*j+1)
        sign = (-1)**j
        term = sign*B
        lo += term//d
        hi += ceildiv(term, d)
        if d > B:
            return lo-1, hi+1
        j += 1

a, b = atan_inv(5)
e, f = atan_inv(239)
pi = 16*a-4*f, 16*b-4*e
root = isqrt(234446*B*B)
c = (2*B*pi[0]//(root+1), ceildiv(2*B*pi[1],root))

def coefficients(limit):
    spf = list(range(limit+1))
    for p in range(2, isqrt(limit)+1):
        if spf[p] == p:
            for j in range(p*p,limit+1,p):
                if spf[j] == j:
                    spf[j] = p
    a = [0]*(limit+1)
    a[1], a[2] = 1, -1
    for n in range(3,limit+1):
        p = spf[n]
        if p == n:
            qr = bytearray(p)
            for j in range(1,(p+1)//2):
                qr[j*j%p]=1
            a[p] = 0
            for x in range(p):
                y = ((x*x-1267)*x+17230)%p
                if y:
                    a[p] += 1-2*qr[y]
        else:
            j = n//p
            a[n] = a[p]*a[j]
            if p != 2 and j%p == 0:
                a[n] -= p*a[j//p]
    return a

def moments():
    a = coefficients(ceildiv(300*B,c[0]))
    out = {2:[0,0], 4:[0,0]}
    for j in range(1,321):
        t = j*B//32
        y = exponential((t,t))
        x = mul(c,y)
        q = exponential((-x[1],-x[0]))
        count = ceildiv(300*B,x[0])
        lo = hi = 0
        for n in range(count,0,-1):
            lo,hi=mul((lo+a[n]*B,hi+a[n]*B),q)
        lo,hi=mul((lo,hi),y)
        for k in out:
            num,den=2*j**k,32**(k+1)
            out[k][0] += num*lo//den
            out[k][1] += ceildiv(num*hi,den)
    return out

if __name__ == "__main__":
    for k, I in moments().items():
        print(k, *(Decimal(z).scaleb(-D) for z in I))
\end{verbatim}
\endgroup
\raggedright
\begin{thebibliography}{9}
\bibitem{curve} LMFDB Collaboration, \emph{Elliptic curve 234446.a1}.
\url{https://www.lmfdb.org/EllipticCurve/Q/234446/a/1}
\bibitem{reliability} LMFDB Collaboration,
\emph{Reliability of elliptic curve data over \(\mathbb Q\)}.
\url{https://www.lmfdb.org/EllipticCurve/Q/reliability}
\bibitem{Wiles} A. Wiles, \emph{The Birch and Swinnerton-Dyer Conjecture}.
\url{https://www.claymath.org/wp-content/uploads/2022/05/birchswin.pdf}
\bibitem{TW} L. N. Trefethen and J. A. C. Weideman,
\emph{The exponentially convergent trapezoidal rule},
SIAM Review 56 (2014), 385--458.
\url{https://people.maths.ox.ac.uk/trefethen/publication/PDF/2014_149.pdf}
\bibitem{decimal} Python Software Foundation, \emph{decimal}.
\url{https://docs.python.org/3/library/decimal.html#decimal.Decimal.exp}
\end{thebibliography}
\end{document}
